% =====================================================================
% READING DOCUMENT (not a submission): current state + diagnosis of the
% norm-boundary project, typeset in the CCAI/NeurIPS workshop style that
% lives in this repo (tackling_climate_workshop_style.sty). Loaded with
% [preprint] so it is de-anonymized and carries no review line numbers.
% Compile: cd paper && tectonic state_diagnosis.tex
% =====================================================================
\documentclass{article}

\PassOptionsToPackage{numbers,compress}{natbib}
\usepackage[preprint]{tackling_climate_workshop_style}

\usepackage[utf8]{inputenc}
\usepackage[T1]{fontenc}
\usepackage{amsmath,amssymb,amsfonts}
\usepackage{graphicx}
\usepackage{booktabs}
\usepackage{nicefrac}
\usepackage{microtype}
\usepackage{xcolor}
\usepackage{hyperref}
\usepackage{url}

\newcommand{\lps}{\ensuremath{\mathrm{LPS}}}
\newcommand{\mse}{\ensuremath{\mathrm{MSE}}}
\newcommand{\var}{\ensuremath{\mathrm{Var}}}
\newcommand{\rawm}{\textsc{Raw}}
\newcommand{\inm}{\textsc{In}}
\newcommand{\cnm}{\textsc{Cn}}
\newcommand{\good}{\textcolor{green!55!black}{\textbf{solid}}}
\newcommand{\bad}{\textcolor{red!75!black}{\textbf{not solid}}}
\newcommand{\defect}{\textcolor{orange!85!black}{\textbf{known defect}}}

\title{norm-boundary: Current State and Diagnosis\\
{\large\normalfont A reader's map of what was run, what came out, and what holds}}

\author{%
  Jaehong Yu, Jaesung Hong, Sungwon Lee\\
  \texttt{internal reading copy --- not for circulation}
}

\begin{document}
\maketitle

\begin{abstract}
This is a reading document, not a paper submission. It records the current
state of the norm-boundary study organized by \emph{what each experiment
tests} rather than by the order it was run, with every headline number
re-derived from the frozen result CSVs. The one-line thesis: the endogenous
instance-normalization layer that is now the default in deep forecasters has
an applicability boundary---when a series' level is predictable from exogenous
covariates (weather for wind and solar, temperature for load), subtracting the
lookback-window mean destroys signal that no downstream capacity restores---and
a pre-training score (\lps) places any series on the correct side of it. Green
marks what is solid; red marks the claims that did not survive verification;
orange marks the known data defect. The exhaustive technical ledger is
\texttt{results/RESULTS.md}.
\end{abstract}

\section{The one thesis, and the evidence in order of strength}

Instance normalization (RevIN and its successors SAN, FAN, Dish-TS) removes each
window's mean and variance before the backbone and restores them at the output,
and every member of the family estimates those statistics \emph{endogenously},
from the target's own past. On exogenously driven energy series this discards a
level the covariates carry. The contribution is not a new normalizer---the arm
we use as an instrument, \cnm, is beaten by cheaper alternatives on the very
series where the rule fires---but the \emph{boundary} of the default layer and a
\emph{pre-training test} for which side of it a series lies on.

The evidence is presented below in decreasing order of strength, which is also
the order in which a paper should present it: structure first, mechanism second,
the head-line dissociation third, the diagnostic last.

\section{What was run}

\begin{table}[h]
\centering\small
\setlength{\tabcolsep}{4pt}
\caption{The experiments, organized by what each one tests. Total: 4{,}202
training runs plus extensions, over eight datasets (four exogenously driven
energy series---Jeju Wind, GEFCom-Wind/Load/Solar---and four standard LTSF
benchmarks).}
\begin{tabular}{clll r}
\toprule
\# & Experiment & Tests & Setup & Runs \\
\midrule
1 & Representational deficiency (Prop.\ 1) & what \inm{} removes from the class & algebra $+$ float64 $+$ Monte Carlo & --- \\
2 & Information parity (Block B) & the layer vs.\ the information & equal covariates to every arm & 1{,}133 \\
3 & Main grid (Block A) & does the sign flip across datasets & 5 norms $\times$ 4 backbones $\times h \times$ 5 seeds & 1{,}794 \\
4 & Variance attribution & where the boundary lives & log-MSE factor decomposition & (from A) \\
5 & \lps{} diagnostic (pre-registered) & can a pre-training scalar predict it & OOS $R^2$, $w{=}96$, $\tau{=}0.3$ & 8 sets \\
6 & Synthetic validation (Gate~1) & do the theory crossings match & $\lambda\times h\times L\times$ seed $\times$ norm & 2{,}640 \\
7 & SOTA replication (Blocks C/D) & is it an adapter artifact & arms injected into TimeXer/iTransformer & 1{,}275 \\
8 & Shrinkage rebuttal (G8) & does shrinking remove the diagnostic & re-trained $\alpha$ recalibration & 204 \\
9 & Probabilistic (Block F) & does the boundary reach intervals & quantile arms, pinball $+$ coverage & 644 \\
\bottomrule
\end{tabular}
\end{table}

\section{What came out --- the solid results}

\paragraph{1. Instance normalization removes the window level irreversibly (\good).}
For \emph{any} backbone $f$, the composed forecast is
$\hat y = \bar y + s\,f\!\left((x-\bar y\mathbf 1)/s,\, g\right)$: additively
separable in the window mean $\bar y$ with unit slope. Driving the real code in
float64 confirms it exactly, $\lvert f(x+c)-f(x)-c\rvert \le 2.7\times10^{-12}$
at $c=-5000$. So \inm{} cannot represent a level response with slope other than
one, nor any level--covariate interaction; \rawm{} can. \textbf{Flexibility does
not fix this}: with an unconstrained backbone \rawm{} reaches the Bayes risk
while \inm{} floors at $[(a-1)^2+\kappa^2\var(g)]\,\var(\bar y)(1-\varrho^2)>0$,
so the \inm{}$-$\rawm{} gap \emph{grows} with capacity. An MLP demonstration at
$a{=}1$ confirms it (RevIN excess $0.085$ vs.\ Raw $0.003$).

\paragraph{2. It is the layer, not the information (\good).}
Give every arm the identical past and future covariates and toggle only the
layer. The isolated cost $\mse(\text{RevIN}{+}\text{cov})-\mse(\rawm{+}\text{cov})$
is positive on four of five backbones (cluster-bootstrap 95\% CIs excluding
zero) and null on the fifth:
\begin{center}\small
\begin{tabular}{lccc}
\toprule
Backbone & mean & 95\% CI & cells $>0$ \\
\midrule
LGBM-Cov & $+0.0191$ & $[+0.0116,+0.0269]$ & 11/11 \\
Linear mixer & $+0.0205$ & $[+0.0056,+0.0368]$ & 8/11 \\
MLP mixer & $+0.0238$ & $[+0.0119,+0.0390]$ & 10/11 \\
SegRNN-Cov & $+0.0112$ & $[+0.0058,+0.0172]$ & 10/11 \\
PatchTST-Cov & $+0.0038$ & $[-0.0091,+0.0167]$ & 7/11 \\
\bottomrule
\end{tabular}
\end{center}
The effect is small ($0.004$--$0.024$) because it is the layer cost alone; the
$+0.841$ head-line gap of Table~\ref{tab:main} measures covariate \emph{access},
a different thing.

\paragraph{3. A strategy-level double dissociation (\good).}
Table~\ref{tab:main} adds benchmark columns and a MASE panel to the main grid.
Covariate-conditioned level models (a first-stage regression alone, a classical
dynamic regression, and \cnm) beat the best instance-normalized arm on
\textbf{11 of 12} exogenous comparisons and \textbf{0 of 12} standard ones
($10/12$ and $0/12$ under MASE). Both exceptions are the ridge dynamic
regression on the nonlinear temperature--load response. The result is about the
strategy, not any one method.

\begin{table}[h]
\centering\small
\setlength{\tabcolsep}{3.5pt}
\caption{Main grid, test MSE (train-fit $z$-scale), backbone/horizon/seed mean.
\textbf{Bold} = row minimum. Benchmarks (left) are deterministic; the three
covariate-conditioned columns (DynReg, FirstStg, \cnm) beat every
instance-normalized column on the exogenous group and lose on the standard one.}
\label{tab:main}
\begin{tabular}{lcccc cccc c}
\toprule
& \multicolumn{4}{c}{Benchmarks} & \multicolumn{4}{c}{Endogenous level} & Exog. \\
\cmidrule(lr){2-5}\cmidrule(lr){6-9}\cmidrule(lr){10-10}
Dataset & S.naive & Clim. & DynReg & FirstStg & Raw & RevIN & SAN & FAN & \cnm \\
\midrule
\multicolumn{10}{l}{\emph{Exogenously driven}}\\
Jeju Wind & 1.627 & 1.032 & \textbf{0.285} & 0.294 & 0.689 & 0.697 & 0.703 & 0.772 & 0.293 \\
GEFCom-Wind & 1.988 & 1.231 & 0.171 & \textbf{0.156} & 1.115 & 1.004 & 0.976 & 1.188 & 0.163 \\
GEFCom-Load & 0.686 & 1.190 & 0.371 & 0.137 & 0.359 & 0.343 & 0.333 & 0.368 & \textbf{0.107} \\
GEFCom-Solar & 0.333 & 0.203 & 0.141 & \textbf{0.057} & 0.183 & 0.184 & 0.212 & 0.183 & 0.058 \\
\midrule
\multicolumn{10}{l}{\emph{Standard LTSF}}\\
ETTh1 & 0.670 & 0.883 & 0.479 & 1.403 & 0.396 & \textbf{0.377} & 0.391 & 0.411 & 2.202 \\
ETTh2 & 0.577 & 3.128 & 0.892 & 3.367 & 0.387 & 0.290 & \textbf{0.287} & 0.325 & 6.604 \\
Electricity & 0.217 & 0.389 & 0.199 & 0.322 & 0.150 & 0.146 & 0.144 & \textbf{0.138} & 0.173 \\
Weather & 0.338 & 0.688 & 0.181 & 0.648 & 0.182 & 0.172 & \textbf{0.167} & 0.175 & 0.727 \\
\bottomrule
\end{tabular}
\end{table}

\paragraph{4. The boundary's location (\good).}
The normalization$\times$dataset interaction explains $47.1\%$ of log-MSE
variance against $0.6\%$ for all backbone terms. That share is carried by the
exogenous arm: restricted to the four endogenous arms the interaction is
$0.43\%$. So normalization is near-immaterial \emph{within} the endogenous
family, and the boundary runs between endogenous and exogenous level handling.

\paragraph{5. A pre-training scalar predicts the side (\good, sign only).}
The official \lps{} value of every dataset and the predicted sign of the
RevIN$-$\cnm{} gap were committed to version control (\texttt{cab17c1}) four
minutes before the first grid run. All eight signs were correct, and the record
is unchanged under MSE, MAE, and MASE.

\begin{center}\small
\begin{tabular}{lc c lc}
\toprule
\multicolumn{2}{c}{Exogenous ($\lps\!\ge\!0.3\Rightarrow$\cnm)} & &
\multicolumn{2}{c}{Standard ($\lps\!<\!0.3\Rightarrow$\inm)}\\
\cmidrule(lr){1-2}\cmidrule(lr){4-5}
GEFCom-Load & 0.894 & & Electricity & 0.283 \\
GEFCom-Solar & 0.875 & & Weather & 0.110 \\
Jeju Wind & 0.745 & & ETTh2 & $-0.205$ \\
GEFCom-Wind & 0.744 & & ETTh1 & $-0.717$ \\
\bottomrule
\end{tabular}
\end{center}

\noindent Synthetic validation (Gate~1, 2{,}640 runs) matches the theory
crossings to within $0.015$, and the pattern replicates on TimeXer and
iTransformer (Blocks C/D, mean RevIN$-$\cnm{} gap $+0.41$).

\section{What did not hold --- the honest column}

\paragraph{\lps{} orders the \emph{side}, not the \emph{magnitude} (\bad).}
Pooled over 23 cells the rank correlation of the gap with \lps{} is $+0.783$,
but that is cluster separation: \emph{within} the 11 exogenous cells it reverses
to $\varrho=-0.750$ ($p=0.008$)---the two highest-\lps{} datasets have the two
smallest gaps. The pre-registered ``magnitude increases with \lps'' auxiliary
prediction is withdrawn. \lps{} is a regime classifier, not a calibrated dial.

\paragraph{The aggregate $2^{-8}$ significance is overstated (\bad).}
Eight datasets are two correlated regimes; granting only that the two regimes
were called correctly gives $p\approx0.25$, not $0.004$. The informative content
is per-dataset, and the threshold plateau $\tau\in[0.30,0.70]$ sits in the empty
gap between the two \lps{} clusters, so the boundary is stress-tested at
effectively one dataset (Electricity).

\paragraph{Shrinkage did not adjudicate, but sizing it needs the diagnostic.}
The rebuttal ``just shrink the level and the ETTh2 catastrophe disappears'' was
pre-registered and run. The loss-minimizing shrink coefficient tracks \lps{}
(best at $\alpha{=}0$ when $\lps{<}0$, at $\alpha{=}\mathrm{clip}(\lps)$ near the
threshold, at $\alpha{=}1$ on the exogenous group), so \emph{choosing} how much
to shrink requires exactly what \lps{} measures; the diagnostic-free textbook
recalibration is worse than plain \rawm{} on all four standard datasets. The
catastrophe framing is partly retracted; the case for a pre-training screen
survives on cost and on the asymmetry (\inm's exogenous failure has no post-hoc
repair, \cnm's does).

\paragraph{Probabilistic intervals under-cover (\bad, kept as caveat).}
\cnm{} has the best pinball loss on all four exogenous datasets, but its
nominal-$80\%$ intervals empirically cover $0.50$--$0.66$ there (RevIN
$0.79$--$0.89$), because first-stage uncertainty is not propagated. This
contradicts any ``tighter reserve margins'' impact story, since reserve sizing
is a quantile decision.

\section{Known defect and limitations}

\paragraph{Lead-matching claim vs.\ implementation (\defect).}
The first stage is fitted once per dataset, outside the horizon loop, so no
horizon-dependent covariate band selection happens. \textbf{Jeju $h{=}48$}
consumes the day-ahead band the code itself labels ``$h\!\le\!24$-valid,'' which
compromises the only real-data test of the horizon prediction; \textbf{GEFCom-Load}
uses realized (not forecast) station temperature. GEFCom-Wind/Solar at
$h{\in}\{96,336\}$ are \emph{not} affected---the competition ships the whole
target-month forecast block. Fix: horizon-legal covariates and re-run the Jeju
$h{=}48$ column ($\sim$100 cells).

\paragraph{Other limitations.}
A single evaluation origin per dataset (no rolling origin); MSE-centric metrics
(MASE now added); a bimodal \lps{} panel with nothing between $0.28$ and $0.74$;
no properly specified classical baseline actually run (only a ridge stand-in,
which itself beats \cnm{} on two datasets); and Jeju Wind is self-curated and
developed alongside the method.

\section{Diagnosis --- where this material can go}

\begin{table}[h]
\centering\small
\caption{Honest venue odds. The writing has been fixed this pass; the ceiling is
now set by the material---incremental novelty (a pre-training normalization
selector already exists, though endogenous and hand-tuned), a dominated method,
and a diagnostic that is a classifier rather than a dial.}
\begin{tabular}{lcc}
\toprule
Target & material now & after one experiment $+$ rigor \\
\midrule
IJF / Applied Energy (top-Q1) & 3--8\% & 30--40\% \\
Energy and AI, Neurocomputing, KBS, EAAI (mid-Q1) & 35--45\% & \textbf{55--70\%} \\
TMLR (quality landing; no JCR quartile) & 40--60\% & 80--85\% \\
CCAI workshop (non-archival) & \textbf{45--50\%} & \textbf{70--75\%} \\
\bottomrule
\end{tabular}
\end{table}

\noindent \textbf{Bottom line.} This is honest, well-executed, mid-Q1 material.
Mid-tier Q1 (energy-AI, applied-AI) is within reach after evaluation-rigor work;
top-Q1 is not, as raw material, because the ceiling is novelty and a dominated
method, which writing cannot lift. The one lever that can raise it is a
\textbf{graded-\lps{} experiment}---disaggregate GEFCom's ten wind zones and
three solar plants (already on disk; RLinear $+$ LightGBM, a few GPU-hours) to
turn the classifier into a possible dose--response dial. It rescues or cleanly
kills the magnitude claim either way. The lowest-cost, highest-probability
near-term output is the CCAI workshop, which does not burn the journal path.

\end{document}
